\chapter{HASIL}
\label{bab:hasil}

\textbf{Bab ini merupakan contoh format pelaporan dan tidak memuat hasil penelitian riil.} Tidak ada skor, jumlah sampel empiris, keberhasilan eksperimen, ataupun perbandingan metode yang dinyatakan telah diperoleh. Tabel berisi nama komponen dan deskripsi keluaran yang perlu disiapkan oleh pengguna. Ganti seluruh materi contoh dengan hasil yang dapat ditelusuri ke data dan prosedur penelitian sendiri.

\section{Cakupan Keluaran}

Pengantar hasil perlu menjelaskan keluaran yang tersedia, hubungan dengan metodologi, dan batas penyajiannya. Jika suatu prosedur tidak menghasilkan keluaran yang direncanakan, kondisi tersebut perlu dinyatakan. Tidak tersedianya hasil adalah informasi yang berbeda dari hasil bernilai nol. Memisahkan keduanya membantu pembaca menilai cakupan jawaban tanpa menganggap seluruh tahap selalu berhasil.

Tabel \ref{tab:contoh-keluaran} menunjukkan format ringkas untuk mengelompokkan komponen. Nama pada tabel merupakan unsur dokumen, bukan daftar temuan penelitian. Bentuk ini dapat diganti dengan komponen yang sesuai dengan bidang pengguna, misalnya kelompok data, tahap analisis, kategori pengamatan, atau keluaran pengukuran.

\begin{table}[htbp]
  \centering
  \caption{Contoh format ringkasan komponen keluaran, tanpa hasil empiris}
  \label{tab:contoh-keluaran}
  \begin{tabular}{p{4.0cm}p{8.6cm}}
    \toprule
    Nama komponen & Deskripsi keluaran \\
    \midrule
    Identitas bahan & Keterangan asal, cakupan, dan batas penggunaan bahan \\
    Prosedur analisis & Keterangan hubungan langkah dengan keluaran yang disajikan \\
    Tabel ringkasan & Susunan informasi yang dapat dibandingkan dalam satu tampilan \\
    Gambar alur & Representasi hubungan tahapan yang dijelaskan dalam narasi \\
    Catatan pemeriksaan & Penjelasan kondisi yang perlu diperbaiki atau diperiksa kembali \\
    \bottomrule
  \end{tabular}
\end{table}

\section{Penyajian Unsur Dokumen}

\subsection{Tabel dan Penjelasan}

Tabel membutuhkan judul yang menjelaskan isi tanpa bergantung sepenuhnya pada paragraf lain. Judul tetap ringkas, sedangkan aturan pembacaan yang lebih panjang dapat diterangkan dalam narasi. Nama kolom harus menunjukkan jenis informasi yang disajikan. Jika tabel penelitian memuat besaran, satuan dan arti nilai kosong perlu disebutkan dengan cara yang konsisten.

Sesudah tabel, penulis dapat menjelaskan hubungan yang paling relevan dengan pertanyaan penelitian. Uraian tidak perlu mengulang seluruh isi sel dalam bentuk paragraf. Penjelasan dipilih berdasarkan peran tabel dalam argumentasi. Dalam contoh ini, Tabel \ref{tab:contoh-keluaran} dipakai untuk menunjukkan bahwa komponen dan deskripsinya mempunyai kolom yang terpisah serta dapat dirujuk dari narasi.

\subsection{Gambar dan Persamaan}

Gambar \ref{fig:contoh-alur} pada Bab \ref{bab:metodologi} memperlihatkan rujukan silang antarbab. Penulis dapat merujuk gambar yang telah diterangkan sebelumnya apabila gambarnya masih relevan. Pengulangan gambar perlu dipertimbangkan berdasarkan kebutuhan pembaca dan panjang pembahasan, bukan dilakukan hanya karena bab berubah.

Persamaan \ref{eq:contoh-rerata} dapat dirujuk untuk menjelaskan definisi suatu besaran. Nilai yang dihasilkan pada penelitian sesungguhnya harus mempunyai satuan, definisi, dan asal penghitungan yang dapat diperiksa. Contoh ini tidak memasukkan nilai ke persamaan karena tidak menyediakan observasi empiris. Penyajian formula dan penyajian hasil pengukuran memiliki peran berbeda.

\section{Contoh Tabel yang Melintasi Halaman}

Tabel \ref{tab:contoh-pemeriksaan-panjang} menggunakan lingkungan tabel panjang agar dapat berpindah halaman dengan judul kolom yang diulang. Keterangan ``(lanjutan)'' membantu pembaca mengenali hubungan bagian tabel pada halaman berikutnya. Kolom pertama menamai komponen dan kolom kedua menjelaskan bahan yang dibutuhkan. Isi tetap berupa petunjuk penulisan, tanpa status pemeriksaan yang diklaim sudah lulus.

\begingroup
\setlength{\tabcolsep}{5pt}
\begin{longtable}{p{3.5cm}p{9.0cm}}
  \caption{Contoh daftar komponen dan deskripsi keluaran untuk tabel panjang}
  \label{tab:contoh-pemeriksaan-panjang}\\
  \toprule
  Nama komponen & Deskripsi keluaran \\
  \midrule
  \endfirsthead
  \caption[]{Contoh daftar komponen dan deskripsi keluaran untuk tabel panjang (lanjutan)}\\
  \toprule
  Nama komponen & Deskripsi keluaran \\
  \midrule
  \endhead
  \midrule
  \multicolumn{2}{r}{Bersambung ke halaman berikutnya} \\
  \endfoot
  \bottomrule
  \endlastfoot
  Judul naskah & Judul yang menjelaskan fokus penelitian dan sesuai dengan judul yang telah disepakati. Penggantian judul dilakukan melalui metadata agar bagian yang memakainya tetap konsisten. \\
  Identitas penulis & Nama dan nomor identitas yang diperiksa terhadap data akademik. Nilai pada contoh merupakan penanda generik yang perlu diganti sebelum naskah digunakan. \\
  Identitas pembimbing & Nama pembimbing sesuai penetapan yang berlaku. Nomor identitas serta pengesahan dicantumkan hanya berdasarkan informasi yang dapat dipertanggungjawabkan. \\
  Persetujuan & Redaksi formulir, ruang tanda tangan, dan tanggal berdasarkan proses persetujuan sebenarnya. Halaman contoh tidak menyatakan bahwa persetujuan telah diberikan. \\
  Orisinalitas & Redaksi pernyataan sesuai ketentuan yang berlaku dan praktik penulisan sebenarnya. Ruang formulir tetap kosong sampai penulis menyiapkan pernyataan yang dapat dipertanggungjawabkan. \\
  Prakata & Uraian singkat mengenai konteks penyusunan dan bantuan yang benar-benar diterima. Nama atau ucapan terima kasih tidak ditambahkan secara fiktif. \\
  Abstrak & Ringkasan masalah, tujuan, metode, hasil, dan batas kesimpulan pada penelitian sesungguhnya. Abstrak contoh menerangkan fungsi template karena hasil empiris tidak disediakan. \\
  Abstract & Ringkasan berbahasa Inggris yang mempunyai makna setara dengan abstrak Indonesia. Istilah, cakupan, dan batas hasil perlu diperiksa pada kedua bahasa. \\
  Kata kunci & Istilah yang membantu pembaca mengenali bidang dan fokus naskah. Istilah dipilih berdasarkan materi sebenarnya dan mengikuti ketentuan jumlah yang berlaku. \\
  Daftar isi & Judul bab dan bagian yang berasal dari struktur sumber. Setelah isi berubah, keluaran diperiksa agar nomor halaman serta hierarki judul mengikuti susunan terbaru. \\
  Daftar tabel & Judul tabel yang konsisten dengan tabel di dalam naskah. Judul yang panjang dapat ditata agar daftar tetap terbaca tanpa menghilangkan makna utamanya. \\
  Daftar gambar & Judul gambar dan lokasi halaman yang sesuai dengan penyajian dalam isi. Gambar perlu dirujuk dan dijelaskan, termasuk apabila berasal dari sumber lain. \\
  Daftar lampiran & Pengenal materi pendukung yang dapat ditemukan kembali oleh pembaca. Materi yang dirujuk sebagai lampiran harus benar-benar terdapat pada bagian akhir naskah. \\
  Pendahuluan & Hubungan masalah, tujuan, manfaat, serta batasan yang menjadi dasar penelitian. Uraian disusun berdasarkan sumber dan kondisi yang dapat diperiksa. \\
  Landasan kepustakaan & Konsep dan penelitian terdahulu yang digunakan dalam argumentasi. Setiap sumber memiliki peran yang dapat dikenali dalam pemilihan metode atau penafsiran hasil. \\
  Metodologi & Prosedur, unit pengamatan, bahan, dan aturan analisis yang menjawab rumusan masalah. Bagian ini menjelaskan asal keluaran yang dilaporkan pada bab hasil. \\
  Hasil & Bukti yang benar-benar diperoleh beserta cakupan dan kondisi pelaporannya. Prosedur yang tidak menghasilkan keluaran dijelaskan tanpa membuat nilai pengganti. \\
  Pembahasan & Hubungan antara bukti, pertanyaan penelitian, konsep, dan keterbatasan. Penafsiran dibatasi oleh cakupan bahan serta prosedur yang digunakan. \\
  Penutup & Jawaban ringkas yang mengikuti bukti dan pembahasan. Saran berhubungan dengan keterbatasan yang telah dijelaskan, bukan tambahan hasil yang belum dianalisis. \\
  Sitasi & Pengenal sumber yang mendukung pernyataan terkait. Penulis memeriksa bahwa kutipan dan parafrasa sesuai dengan isi serta batas sumber asal. \\
  Daftar referensi & Identitas bibliografis sumber yang dipakai dalam naskah. Jenis sumber, penulis, judul, dan alamat diperiksa tanpa mengisi data penerbitan dengan nilai rekaan. \\
  Persamaan & Definisi simbol, syarat penggunaan, serta rujukan pada narasi. Persamaan yang memerlukan beberapa baris ditata dengan nomor yang sesuai kebutuhan rujukan. \\
  Lampiran & Materi pendukung yang terlalu rinci untuk alur utama tetapi diperlukan untuk pemeriksaan. Bahan yang dibatasi aksesnya tidak disertakan tanpa dasar izin yang sesuai. \\
  Keluaran akhir & Dokumen yang diperiksa kembali sesudah perubahan terakhir. Pemeriksaan mencakup tampilan, rujukan, identitas, serta kelengkapan yang diperlukan untuk penggunaan akademik. \\
\end{longtable}
\endgroup

\section{Kelengkapan Pelaporan}

Pelaporan perlu mempertahankan hubungan antara bahan, prosedur, dan keluaran. Apabila informasi disingkat dalam tabel utama, rincian dapat diletakkan pada lampiran dan dirujuk secara jelas. Informasi penting tentang batas bukti tidak sebaiknya hanya diletakkan di berkas terpisah yang tidak disebut dalam naskah.

Pada contoh ini, Lampiran \ref{lamp:contoh} menyediakan susunan catatan pendukung yang generik. Pengguna dapat mengganti catatan tersebut dengan bahan yang relevan dan berizin. Bab \ref{bab:pembahasan} memperlihatkan cara meneruskan pelaporan menuju penafsiran, tanpa menjadikan contoh format sebagai klaim bahwa hasil ilmiah telah tersedia.
